% !TEX root = ../main_mk.tex
\chapter{Вовед}\label{ch:intro}

\section{Мотивација}

Стилот на еден автор е она што останува од песната кога ќе се смени нејзината
тема: зборовите по кои посегнува авторот, должината на стиховите, дали
има рима и каква, каква слика остава. Не е едноставно да се препише една песна во стил на друг автор, а патем да се задржи содржината. Ова е тежок тест за дали јазичниот модел го научил тој слој на јазикот. Ова е задачата на \emph{пренос на стил на текст} (text style transfer), во најтежок облик: „стилот“ не е нешто објективно, како формален или неформален говор. Не секогаш има еден точен одговор.

Најголемиот дел од истражувањата за пренос на стил, и за јазичните модели
воопшто, се прават на англиски јазик, каде што податоците се изобилни.
Македонскиот јазик е со ограничени ресурси: има малку анотирани корпуси,
малку алатки, и сочинува мал дел од текстовите на кои се обучуваат
општите модели. Поезијата ја влошува оваа состојба. Еден поет во текот на
животот пишува неколку стотини песни, често и многу помалку, па дури и
целокупниот корпус на еден автор е мало податочно множество според мерилата
на машинско учење. Корпусот на врз кој е изградена оваа работа содржи 1.239 
песни и текстови на песни од 38 македонски автори; најзастапениот автор
има 182 песни, следните двајца 160 и 134, а половина од авторите имаат
седумнаесет или помалку.

Големите модели можат со соодветно барање (prompt) да пишуваат во даден
стил, но нивното извршување е скапо и не е под контрола на корисникот.
Малите модели, со неколку милијарди параметри, можат да се обучуваат и
извршуваат на една графичка картичка, што ги прави практичен избор за
јазични задачи без големи пресметковни ресурси. Дали можат да научат
задача толку специфична како гласот на одреден поет, и тоа од толку малку
податоци, е прашањето на кое оваа работа се обидува да одговори.

\section{Истражувачко прашање}

Работата го поставува следното прашање:

\begin{quote}
\emph{Може ли мал јазичен модел, кому му се дадени само изворна песна и
името на македонски автор, да ја препише песната во стилот на тој автор,
и дали учењето со поттикнување го подобрува резултатот?}
\end{quote}

Два дела од прашањето бараат појаснување. „Во стилот на тој автор“ се мери
на два меѓусебно независни начини: автоматски, преку стилометриски
карактеристики и преку класификатор на автори што обуката никогаш не
го гледа, и преку човечка проценка. „Учење со поттикнување“ (reinforcement
learning) е методот Group Relative Policy Optimization (GRPO)\footnote{Според Правописот на македонскиот
литературен јазик (2017), називите од туѓо потекло се предаваат со кирилско писмо,
но се допуштаат и латинични скратеници. Во оваа работа називите на моделите,
методите и алатките (Qwen, Gemma, LoRA, GRPO и сл.) се пишуваат оригинално, на
латиница, за да се избегне обременување на текстот и за да се препознаваат при
пребарување литература од областа.}
\citep{shao2024deepseekmath}, применет по надгледуваното фино
подесување, со награда изградена од стилометрискиот профил на секој автор.

\section{Пристап и главен наод}

Работата го доби својот конечен облик по неколку пристапи што не успеаја.
Замената на избрани зборови од изворната песна со зборови карактеристични
за целниот автор, соодветно менувани по форма, ја зачувува песната, но
малку го менува нејзиниот стил. Кога од мали локални модели се бара да ги
пополнат таквите празни места, тие не успеваат веќе на ниво на
македонскиот јазик. Она што успеа е голем модел, Gemma~4 31B
\citep{gemma4}, на кој му е дадено детално структурирано барање: резиме на
изворната песна, опис на стилот на целниот автор во проза, палета од
зборови и примерни стихови, и експлицитни структурни цели. Тој модел стана
\emph{учител}. Неговите песни, околу 115 за секој од десетте целни автори,
се податоците за обука на двајца \emph{ученици} со слична големина,
Qwen3-4B \citep{yang2025qwen3} и Gemma~4 E4B, кои ги гледаат само изворната песна
и името на авторот. Секој ученик најпрвин се фино подесува со LoRA
\citep{hu2022lora} врз песните на учителот, а потоа
понатаму се обучува со GRPO.

Главниот наод е негативен. Учениците ја научуваат формата на задачата:
пишуваат на македонски, пишуваат песни, а не проза, и ја менуваат
изворната песна наместо да ја препишат. Гласовите не ги научуваат.
Независен класификатор на автори го препознава бараниот автор кај 16,7\%
од песните на фино подесениот ученик Qwen на тест-паровите, наспроти 66,3\% кај
вистински, невидени песни од истите автори; по GRPO оваа бројка паѓа на
12,4\%, иако наградата што методот ја оптимизира расте. Учителот, пак,
со целиот опис на авторот пред себе, го постигнува тоа кај 53,6\% од
своите песни. \pending{Резултати за ученикот Gemma и човечката проценка.}
Препознаените песни на учениците речиси сите им припаѓаат на трите
автори со најмногу песни во корпусот, додека учителот успешно имитира и
некои автори со помалку од четириесет песни.

Работата тврди дека ограничувањето е во податоците, а не во методот.
Корпусот е доволен за учителот да имитира автор кога доказите за неговиот
стил му се дадени во секое барање; не е доволен ученикот да го научи тој
стил и да го задржи во своите параметри од околу 115 примери по автор.
Посилен ученик, подобро подесена награда и учењето со поттикнување не ја
надминуваат оваа граница.

\section{Придонеси}

\begin{itemize}
  \item Исчистен корпус на македонска поезија, означен со зборовни
        видови, со идентификатори на ниво на песна, документирани
        корекции на означувањето и фиксна поделба по автор на
        множества за обука, валидација и тестирање.
  \item Стилометриски профили на авторите и награда за пренос на стил
        изградена од нив: распределби на зборовни видови, лексичка
        разновидност, структура на стихови и строфи, рима,
        карактеристичен речник и класификатор на автори, со порти
        (gates) против препишување на изворот и против дегенериран излез.
  \item Систем учител–ученик што го дестилира големиот модел со
        структурирано барање во двајца мали ученици, по што следи
        GRPO, со исти поставки за двајцата ученици за да можат нивните
        резултати да се споредат.
  \item Евалуација што ги одделува вистинските подобрувања од
        „измамата“ на наградата: класификаторот што се користи како
        награда никогаш не се користи за проценка на резултатот; тоа го
        прави втор класификатор, обучен врз други карактеристики.
  \item Човечка проценка на учителот, двајцата ученици и вистински
        песни врз истите тест-парови, со вистинските песни како слепа
        контрола.
\end{itemize}

\section{Структура на работата}

Поглавјето~\ref{ch:background} дава преглед на преносот на стил на текст,
стилометријата, дестилацијата на знаење и учењето со поттикнување кај
јазичните модели. Поглавјето~\ref{ch:data} го опишува корпусот и неговата
подготовка, а Поглавјето~\ref{ch:stylometry} стилометриските профили и
класификаторите врз кои се потпира остатокот од работата.
Поглавјето~\ref{ch:path} ги сумира претходните пристапи и причините зошто
беа напуштени. Поглавјето~\ref{ch:method} ги претставува учителот,
учениците и наградата, а Поглавјето~\ref{ch:eval} протоколот за
евалуација. Поглавјето~\ref{ch:results} ги дава резултатите,
Поглавјето~\ref{ch:discussion} дискутира зошто учениците заостануваат, а
Поглавјето~\ref{ch:conclusion} ја заклучува работата.
